This appendix consolidates methodological detail moved from Chapter~\ref{ch:method}
to keep the main methodology chapter readable for a time-limited jury review.

\section{Configuration, Provenance, and Ethics}
\label{sec:app-method-config}

Configuration surfaces include theme weights for use-case readiness
(Equation~\eqref{eq:usecase}), filter thresholds in the UI, and embedding model
choice for future NLP. These should be logged, justified, and held fixed when
comparing dashboard releases.

Every country record stores legislation titles and narrative fields traceable to
agency documents or peer-reviewed analyses. When secondary sources disagree, the
protocol prefers: (1)~primary legal text; (2)~official agency guidance;
(3)~peer-reviewed comparative studies; (4)~reputable policy observatories.

Publishing comparative scores can create reputational effects for countries.
Scores are research instruments, not rankings for investment or diplomacy.
Emerging jurisdictions may score lower because documentation is scarce in English
or because AI-SaMD pathways are genuinely early---two different situations that
narratives must disentangle.

New dataset contributors should read the rubric table, shadow one country update
end-to-end, propose scores with written justifications, and merge only after
supervisor or peer review.

\section{Ingestion, Preprocessing, and Infrastructure}
\label{sec:app-method-ingest}

Sources enter the pipeline through a lightweight ETL pattern. Agency PDFs and HTML
guidance pages are logged with URL, access date, and jurisdiction tags.
Bibliographic hits are exported to BibTeX and deduplicated by DOI/title. Country
records are updated in JSON rather than overwritten silently.

For the planned NLP scale-up, documents would be converted to plain text,
segmented into clauses, normalized, and embedded with a transformer
encoder~\cite{vaswani2017,devlin2019,beltagy2019}. Train/validation/test splits
should be stratified by jurisdiction and theme. Infrastructure for the delivered
MVP is intentionally minimal: a local Python environment, Streamlit for UI, and
MiKTeX/TeX Live for report compilation.

\section{Dashboard Specification, Evaluation, and Runbook}
\label{sec:app-method-funcspec}

The information architecture follows progressive disclosure: world overview,
filtered comparison, country narratives, and literature/methodological caveats.
The Streamlit application implements six modules:
\begin{enumerate}[noitemsep]
  \item \textbf{World Map:} choropleth of composite or selected theme scores;
  \item \textbf{Country Comparison:} multi-select radar and bar charts;
  \item \textbf{Theme Analysis:} heatmap, regional means, and gap charts;
  \item \textbf{Global Trends:} adoption intensity and emergence year cards;
  \item \textbf{Country Details:} narrative accordion fields;
  \item \textbf{Literature Review:} in-app markdown rendering.
\end{enumerate}
Filters for region, maturity, and theme thresholds apply globally where relevant.
CSV export preserves auditability outside the UI.

Because the primary artifact is curated comparative analytics, evaluation uses
research-software metrics rather than clinical AUC alone: coverage (20$\times$7
theme matrix), schema consistency, clean-environment dashboard launch, comparative
utility for theme-gap analysis, and planned NLP precision/recall/F1 once labels
mature.

A third party should be able to: (1)~clone the GitHub repository;
(2)~create a virtual environment and install \texttt{requirements.txt};
(3)~run \texttt{streamlit run app.py}; (4)~open the dataset JSON and verify theme
keys; (5)~regenerate report figures via \texttt{generate\_figures.py};
(6)~compile the \LaTeX\ sources. Before any public release, confirm all theme
scores lie in $1..10$, bibliography keys resolve, and the report compiles without
undefined references.
